\documentclass[10pt, aspectratio=169]{beamer}
\input{../../_en_preambulo_beamer}
% Comandos y macros estadísticas compartidas para las presentaciones Beamer en inglés (Metropolis)

% Conjuntos numéricos básicos
\newcommand{\R}{\mathbb{R}}
\newcommand{\N}{\mathbb{N}}
\newcommand{\Q}{\mathbb{Q}}
\newcommand{\Z}{\mathbb{Z}}
\newcommand{\set}[1]{\left\{ #1 \right\}}
\newcommand{\abs}[1]{\left|#1\right|}
\newcommand{\norm}[1]{\left\| #1 \right\|}

% Comandos estadísticos y probabilísticos del curso
\newcommand{\Var}{\operatorname{Var}}
\newcommand{\cov}{\operatorname{Cov}}
\newcommand{\corr}{\rho}
\newcommand{\comb}[2]{\begin{pmatrix} #1 \\ #2 \end{pmatrix}}
\newcommand{\s}{\sigma}
\newcommand{\particion}{\mathfrak{P}}
\newcommand{\card}[1]{\#\left( #1 \right)}
\newcommand{\seq}[3]{\set{#1}_{#2}^{#3}}
\newcommand{\E}{\mathbb{E}}
\newcommand{\Prob}{\mathbb{P}}

% Entornos pedagógicos y cajas en inglés académico natural (soportando tanto nombres en inglés como en español)
\newenvironment{discussion}{%
  \begin{alertblock}{Discussion Question}%
}{%
  \end{alertblock}%
}
\newenvironment{discusion}{%
  \begin{alertblock}{Discussion Question}%
}{%
  \end{alertblock}%
}

\newenvironment{reflection}{%
  \begin{alertblock}{Classroom Reflection}%
}{%
  \end{alertblock}%
}
\newenvironment{reflexion}{%
  \begin{alertblock}{Classroom Reflection}%
}{%
  \end{alertblock}%
}

\newenvironment{paradox}{%
  \begin{alertblock}{Exploratory Paradox}%
}{%
  \end{alertblock}%
}
\newenvironment{paradoja}{%
  \begin{alertblock}{Exploratory Paradox}%
}{%
  \end{alertblock}%
}

\newenvironment{motivation}{%
  \begin{exampleblock}{Motivation: Data Science in Practice}%
}{%
  \end{exampleblock}%
}
\newenvironment{motivacion}{%
  \begin{exampleblock}{Motivation: Data Science in Practice}%
}{%
  \end{exampleblock}%
}

\newenvironment{quickchallenge}{%
  \begin{block}{Quick Team Challenge}%
}{%
  \end{block}%
}
\newenvironment{retoclase}{%
  \begin{block}{Quick Team Challenge}%
}{%
  \end{block}%
}


\title[Model Validation]{Section 09.08: Model Validation}
\subtitle{Train/Test Split and $k$-Fold Cross-Validation}
\author[J. Castillo Colmenares]{Juliho Castillo Colmenares}
\institute{Tecnológico de Monterrey}
\date{\vspace{-1.5cm}}

\begin{document}

% Slide 1: Title
\begin{frame}[plain]
  \titlepage
\end{frame}

% Slide 2: Roadmap
\begin{frame}{Roadmap --- Chapter 09: Linear and Multiple Regressions}
  \scriptsize
  Where are we in the modular development of this chapter?
  \begin{itemize}\itemsep=0.02em
    \item \textbf{09.01} Correlation as the Premise of Regression
    \item \textbf{09.02} Introduction to Linear Regression
    \item \textbf{09.03} The Mathematics of Regression: OLS Derivation and $R^2$
    \item \textbf{09.04} Linear Regression on Simulated Data
    \item \textbf{09.05} Optimal Coefficients, $t$/$F$ Tests, and RSE
    \item \textbf{09.06} Implementation in Python with \texttt{statsmodels}
    \item \textbf{09.07} Multiple Regression, the Normal Equation, and Ridge/Lasso
    \item \textbf{09.08} Model Validation and $k$-fold Cross-Validation \emph{(Today)}
    \item \textbf{09.09} Residual Diagnostics and Multicollinearity (VIF)
    \item \textbf{09.10} Regression with \texttt{scikit-learn}
    \item \textbf{09.11} Categorical Variables and Dummy Variables
    \item \textbf{09.12} Nonlinear Transformations and Polynomial Regression
  \end{itemize}
  \pause
  \vspace{-0.05cm}
  \begin{block}{Session Objective}
    Determine whether a fitted model \textbf{generalizes} to new data: implement the train/test split to detect overfitting, build $k$-fold cross-validation from scratch, and use it to robustly compare competing models.
  \end{block}
\end{frame}

% Slide 3: Motivation
\begin{frame}{Does the Model Generalize, or Just Memorize?}
  \footnotesize
  \begin{motivacion}
    A predictive model needs to be validated by observing its performance on \textbf{data different} from what was used to fit it. A model can fit the training data extremely well (\emph{overfitting}) and still completely fail to predict new observations.
  \end{motivacion}

  \vspace{0.15cm}
  \pause
  All diagnostics from Sections 09.03-09.07 (including $R^2$, $t$/$F$ tests) are computed \textbf{within} the fitting sample: none of them, by itself, guarantees good performance outside of it.
\end{frame}

% Slide 4: Train/test split
\begin{frame}{Train/Test Split}
  \footnotesize
  \begin{block}{Procedure}
    We randomly split the dataset (typically 80\%/20\%): we fit the model \textbf{only} with the training data, and evaluate its error \textbf{only} with the test data, never seen during fitting.
  \end{block}
  \pause

  \vspace{0.1cm}
  \begin{alertblock}{Overfitting Signal}
    If the error (RSE/RMSE) on the test set is substantially larger than on the training set, the model memorized peculiarities of the fitting sample instead of capturing the true underlying relationship.
  \end{alertblock}
\end{frame}

% Slide 5: k-fold cross-validation
\begin{frame}{$k$-Fold Cross-Validation}
  \footnotesize
  \begin{block}{Procedure}
    We split the data into $k$ subsets (folds) of approximately equal size. In each of the $k$ iterations: we train on $k-1$ folds and evaluate on the remaining fold. We average the $k$ metrics obtained.
  \end{block}
  \pause

  \vspace{0.1cm}
  \begin{alertblock}{Advantages over a Single Split}
    Lower variance (any single particular partition influences the final estimate less); efficient use of data (each observation is used exactly once for testing); detects overfitting more robustly than a single split.
  \end{alertblock}
\end{frame}

% Slide 6: Choosing k
\begin{frame}{Choosing $k$ and Stratified Cross-Validation}
  \footnotesize
  \begin{itemize}\itemsep=0.1cm
    \item $k=5$ or $k=10$: conventional values, good bias-variance balance. \pause
    \item Small $k$: higher bias, lower variance, faster. \pause
    \item $k=n$ (leave-one-out, LOOCV): nearly unbiased, but slow and high-variance. \pause
    \item \textbf{Stratified CV:} preserves the distribution of a variable of interest across folds, useful with imbalanced data.
  \end{itemize}
\end{frame}

% Slide 7: Python Lab (Block 1)
\begin{frame}[fragile]{Python Lab: Detecting Overfitting}
  \scriptsize
  An overparameterized model (16 predictors for $n=40$) fits training data well but fails on test (\texttt{09.08\_model\_validation.py}):
  {\lstset{basicstyle=\fontsize{3.0pt}{3.7pt}\selectfont\ttfamily,aboveskip=0.05em,belowskip=0.05em}
  \lstinputlisting[firstline=36, lastline=62]{../../code/09_regresiones/09.08_model_validation.py}}
\end{frame}

% Slide 8: Python Lab (Block 2)
\begin{frame}[fragile]{Python Lab: $k$-Fold from Scratch}
  \scriptsize
  Full implementation of 5-fold cross-validation with no external libraries:
  {\lstset{basicstyle=\fontsize{2.9pt}{3.5pt}\selectfont\ttfamily,aboveskip=0.05em,belowskip=0.05em}
  \lstinputlisting[firstline=65, lastline=85]{../../code/09_regresiones/09.08_model_validation.py}}
\end{frame}

% Slide 9: Python Lab (Block 3)
\begin{frame}[fragile]{Python Lab: Comparing Models via CV}
  \scriptsize
  Robust comparison between a single-predictor model and a two-predictor model:
  {\lstset{basicstyle=\fontsize{3.2pt}{3.9pt}\selectfont\ttfamily,aboveskip=0.05em,belowskip=0.05em}
  \lstinputlisting[firstline=106, lastline=119]{../../code/09_regresiones/09.08_model_validation.py}}
\end{frame}

% Slide 10: Terminal Output
\begin{frame}[fragile]{Python Lab: Terminal Output}
  \scriptsize
  Standard output produced when running the lab script:
  \vspace{0.02cm}
  \begin{block}{Standard Console Output}
    \fontsize{3.4pt}{4.1pt}\selectfont
    \begin{verbatim}
=== Block 1: Overfitting Detection ===
n_train=32, n_test=8, p=16 (1 real + 15 pure noise)
RSE train=8.0786, RMSE test=12.4452 (diff=4.3666) -> overfitting signature

=== Block 2: k-Fold CV from Scratch ===
Folds: 0.9056, 0.8083, 0.8765, 0.8452, 0.8409
Mean R^2 = 0.8553 +/- 0.0332

=== Block 3: Comparing Models via CV ===
Model 1 (x1 only):  mean R^2 = 0.5727 +/- 0.0804
Model 2 (x1+x2):    mean R^2 = 0.8553 +/- 0.0332
    \end{verbatim}
  \end{block}
\end{frame}

% Slide 11: Interpretation
\begin{frame}{Interpretation: Why Does Validation Matter?}
  \footnotesize
  \begin{itemize}\itemsep=0.12cm
    \item Block 1 shows the danger of having \textbf{more predictors than genuine information}: 15 pure-noise columns artificially reduce the training error, but the test error exposes it. \pause
    \item Block 3 confirms, with \emph{out-of-sample} evidence, that the second genuine predictor substantially improves the model (average $R^2$ from $0.57$ to $0.86$) --- a more reliable conclusion than comparing in-sample $R^2$. \pause
    \item The standard deviation across folds ($\pm0.033$ vs. $\pm0.080$) also reveals that the full model is more \textbf{stable}, not just more accurate on average.
  \end{itemize}
\end{frame}

% Slide 12: Closing
\begin{frame}{Section Closing: Model Validation}
  \begin{reflexion}
    No diagnostic computed within the fitting sample --- neither $R^2$ nor the $t$/$F$ tests --- guarantees that a model generalizes. $k$-fold cross-validation is the industry's robust standard for honestly estimating a model's performance on data it has never seen.
  \end{reflexion}

  \vspace{0.2cm}
  \pause
  \begin{block}{Key Takeaways}
    \begin{itemize}\itemsep=0.08cm
      \item \textbf{Train/test split:} fast but high-variance; useful for initial exploration.
      \item \textbf{$k$-fold CV:} averages over $k$ partitions, reducing the variance of the error estimate.
      \item \textbf{Model comparison:} CV is more reliable than comparing in-sample $R^2$.
    \end{itemize}
  \end{block}
\end{frame}

% Slide 13: Modular Perspective
\begin{frame}{Modular Perspective --- The Next Challenge}
  \scriptsize
  \begin{retoclase}
    We have validated the model \textbf{externally} (generalization to new data). Section 09.09 audits it \textbf{internally}: we will verify the classical assumptions about the errors (normality, homoscedasticity, independence) and quantify multicollinearity via the Variance Inflation Factor --- two entirely different, complementary forms of trustworthiness.
  \end{retoclase}

  \vspace{0.08cm}
  \pause
  \begin{alertblock}{Recommended Preparation}
    \begin{itemize}\itemsep=0.06cm
      \item Reflect on why a model can generalize reasonably well (good CV performance) and still violate the classical regression assumptions.
      \item Review the Hat Matrix from Section 09.07: its diagonal will be key to diagnosing influential observations.
    \end{itemize}
  \end{alertblock}
\end{frame}

\end{document}
